\section{Kommunikation und Schnittstellen}
\label{sec:kommunikation}

\subsection{Kommunikationsprotokoll}

Alle Kommunikation zwischen den Komponenten von TaskGrid+ erfolgt
ausschließlich über \textbf{gRPC} (Google Remote Procedure Call)
mit \textbf{Protocol Buffers (Protobuf)} als Serialisierungsformat.
gRPC wurde aus folgenden Gründen gewählt:

\begin{itemize}
  \item Effiziente binäre Serialisierung: deutlich kompakter als JSON/REST.
  \item Strenge Typisierung durch das \texttt{.proto}-Schema: verhindert
        Schnittstellenfehler zur Laufzeit.
  \item Bidirektionale Kommunikation möglich (hier nicht genutzt, aber vorbereitet).
  \item Sprachunabhängigkeit: Client, Dispatcher, Worker und Namensdienst
        können in verschiedenen Sprachen implementiert sein.
\end{itemize}

Die Schnittstellendefinition liegt in \texttt{proto/taskgrid.proto}.

\subsection{Protobuf-Nachrichten}
Alle Nachrichten folgen einem einheitlichen Schema:

\begin{itemize}
	\item \texttt{message\_type}: Typ der Nachricht
	\item \texttt{request\_id}: eindeutige Korrelations-ID
	\item \texttt{timestamp}: Erstellungszeitpunkt der Nachricht
	\item \texttt{sender}: Absender der Nachricht
	\item \texttt{payload}: nachrichtenspezifische Nutzdaten
\end{itemize}
Die Felder \texttt{message\_type}, \texttt{request\_id},
\texttt{timestamp} und \texttt{sender} sind in allen Nachrichten identisch.
Lediglich der Inhalt des Payloads unterscheidet sich je nach Nachrichtentyp.

\subsubsection*{Task-Nachrichten (Dispatcher $\leftrightarrow$ Worker)}

\begin{table}[H]
	\centering
	\begin{tabular}{lll}
		\toprule
		\textbf{Nachricht} & \textbf{Payload} & \textbf{Bedeutung} \\
		\midrule
		\multirow{4}{*}{\texttt{TaskRequest}} 
		& \texttt{task\_id: int32} & Eindeutige Task-Nummer \\
		& \texttt{task\_type: string} & Aufgabentyp \\
		& \texttt{task\_payload: string} & Eingabedaten der Aufgabe \\
		& \texttt{task\_status: string} & Aktueller Status der Aufgabe \\
		\midrule
		\multirow{5}{*}{\texttt{TaskResponse}} 
		& \texttt{accepted: bool} & Gibt an, ob der Worker die Aufgabe annimmt \\
		& \texttt{task\_id: int32} & Gespiegelte Task-ID \\
		& \texttt{worker\_id: string} & Bearbeitender Worker \\
		& \texttt{status: string} & Status der Annahme oder Bearbeitung \\
		& \texttt{error: string} & Fehlermeldung, falls vorhanden \\
		\bottomrule
	\end{tabular}
	\caption{Task-Nachrichten zwischen Dispatcher und Worker}
\end{table}

\subsubsection*{Ergebnis-Rückmeldung (Worker $\rightarrow$ Dispatcher)}

\begin{table}[H]
	\centering
	\begin{tabular}{lll}
		\toprule
		\textbf{Nachricht} & \textbf{Payload} & \textbf{Bedeutung} \\
		\midrule
		\multirow{5}{*}{\texttt{ResultRequest}} 
		& \texttt{task\_id: int32} & Eindeutige Task-Nummer \\
		& \texttt{worker\_id: string} & Worker, der die Aufgabe bearbeitet hat \\
		& \texttt{status: string} & Ergebnisstatus der Aufgabe \\
		& \texttt{result: string} & Ergebnis der Verarbeitung \\
		& \texttt{error: string} & Fehlermeldung, falls vorhanden \\
		\midrule
		\multirow{2}{*}{\texttt{Ack}} 
		& \texttt{success: bool} & Angenommen / Abgelehnt \\
		& \texttt{message: string} & Bestätigungs- oder Fehlermeldung \\
		\bottomrule
	\end{tabular}
	\caption{Ergebnis-Nachrichten vom Worker an den Dispatcher}
\end{table}

\subsubsection*{Client-Schnittstelle (Client $\rightarrow$ Dispatcher)}

\begin{table}[H]
	\centering
	\begin{tabular}{m{4cm} m{5cm} m{6.5cm}}
		\toprule
		\textbf{Nachricht} & \textbf{Feld} & \textbf{Bedeutung} \\
		\midrule
		\multirow{2}{*}{\texttt{PostTaskRequest}} 
		& \texttt{task\_type: string} & Aufgabentyp \\
		& \texttt{task\_payload: string} & Eingabedaten der Aufgabe \\
		\midrule
		\multirow{4}{*}{\texttt{PostTaskResponse}} 
		& \texttt{success: bool} & Angenommen / Abgelehnt \\
		& \texttt{task\_id: int32} & Zugewiesene Task-ID \\
		& \texttt{status: string} & Status der Aufgabe \\
		& \texttt{message: string} & Zusatzinformation zur Anfrage \\
		\midrule
		\multirow{1}{*}{\texttt{GetResultRequest}} 
		& \texttt{task\_id: int32} & ID der abzufragenden Aufgabe \\
		\midrule
		\multirow{5}{*}{\texttt{ResultResponse}} 
		& \texttt{found: bool} & Ergebnis gefunden oder nicht \\
		& \texttt{task\_id: int32} & ID der abgefragten Aufgabe \\
		& \texttt{status: string} & Status der Aufgabe \\
		& \texttt{result: string} & Ergebnis der Aufgabe \\
		& \texttt{error: string} & Fehlermeldung, falls vorhanden \\
		\bottomrule
	\end{tabular}
	\caption{Nachrichten zwischen Client und Dispatcher}
\end{table}

\subsubsection*{Namensdienst-Nachrichten}
\begin{table}[H]
	\centering
	\begin{tabular}{m{4.75cm} m{4.5cm} m{6.5cm}}
		\toprule
		\textbf{Nachricht} & \textbf{Feld} & \textbf{Bedeutung} \\
		\midrule
		\multirow{6}{*}{\texttt{RegisterWorkerRequest}} 
		& \texttt{worker\_id: string} & Eindeutige Worker-ID \\
		& \texttt{task\_types: repeated string} & Unterstützte Aufgabentypen \\
		& \texttt{address: string} & Netzwerkadresse des Workers \\
		& \texttt{port: int32} & Port des Workers \\
		& \texttt{status: string} & Aktueller Worker-Status \\
		& \texttt{current\_load: int32} & Aktuelle Auslastung des Workers \\
		\midrule
		\multirow{3}{*}{\texttt{HeartbeatRequest}} 
		& \texttt{worker\_id: string} & Worker, der den Heartbeat sendet \\
		& \texttt{current\_load: int32} & Aktuelle Auslastung des Workers \\
		& \texttt{status: string} & Aktueller Worker-Status \\
		\midrule
		\multirow{1}{*}{\texttt{DeregisterWorkerRequest}} 
		& \texttt{worker\_id: string} & Zu deregistrierender Worker \\
		\midrule
		\multirow{1}{*}{\texttt{LookupRequest}} 
		& \texttt{task\_type: string} & Gesuchter Aufgabentyp \\
		\midrule
		\multirow{3}{*}{\texttt{LookupResponse}} 
		& \texttt{found: bool} & Worker gefunden oder nicht \\
		& \texttt{workers: repeated WorkerPayload} & Liste passender Worker \\
		& \texttt{message: string} & Zusatzinformation zur Suche \\
		\midrule
		\multirow{2}{*}{\texttt{Ack}} 
		& \texttt{success: bool} & Gibt an, ob die Anfrage erfolgreich war \\
		& \texttt{message: string} & Bestätigungs- oder Fehlermeldung \\
		\bottomrule
	\end{tabular}
	\caption{Namensdienst-Nachrichten}
\end{table}

\subsection{gRPC-Dienste}

Das System definiert drei gRPC-Services in \texttt{proto/taskgrid.proto}:

\subsubsection*{NamingService}

Der Namensdienst implementiert \texttt{NamingService}.
Worker rufen die vier Methoden direkt auf:
\begin{lstlisting}[language={},caption={NamingService in taskgrid.proto}]
service NamingService {
  rpc RegisterWorker   (WorkerInfo)       returns (Ack);
  rpc SendHeartbeat    (HeartbeatRequest) returns (Ack);
  rpc DeregisterWorker (WorkerInfo)       returns (Ack);
  rpc LookupWorker     (LookupRequest)    returns (LookupResponse);
}
\end{lstlisting}

\subsubsection*{DispatcherService}

Der Dispatcher implementiert \texttt{DispatcherService}.
Clients und Worker kommunizieren hierüber:

\begin{lstlisting}[language={},caption={DispatcherService in taskgrid.proto}]
service DispatcherService {
  rpc PostTask     (PostTaskRequest)  returns (TaskResponse);
  rpc GetResult    (GetResultRequest) returns (ResultResponse);
  rpc GetStatus    (StatusRequest)    returns (StatusResponse);
  rpc ReturnResult (ResultRequest)    returns (Ack);
}
\end{lstlisting}

\subsubsection*{WorkerService}

Jeder Worker implementiert \texttt{WorkerService}.
Der Dispatcher ruft \texttt{ExecuteTask} auf:

\begin{lstlisting}[language={},caption={WorkerService in taskgrid.proto}]
service WorkerService {
  rpc ExecuteTask (TaskRequest) returns (TaskResponse);}
\end{lstlisting}

\subsection{Kommunikationsfluss}

Der vollständige Fluss für eine erfolgreiche Task-Verarbeitung:

\begin{enumerate}
  \item \textbf{Client $\rightarrow$ Dispatcher} \texttt{PostTask}: \\
        Client sendet Task-Typ und Payload. Dispatcher gibt \texttt{task\_id} zurück.
  \item \textbf{Dispatcher $\rightarrow$ Namensdienst} \texttt{LookupWorker}: \\
        Dispatcher fragt nach aktiven Workern für den Task-Typ.
  \item \textbf{Dispatcher $\rightarrow$ Worker} \texttt{ExecuteTask}: \\
        Dispatcher sendet Task an ausgewählten Worker per Round-Robin.
  \item \textbf{Worker $\rightarrow$ Dispatcher} \texttt{ReturnResult}: \\
        Worker sendet Ergebnis zurück (COMPLETED oder FAILED).
  \item \textbf{Client $\rightarrow$ Dispatcher} \texttt{GetResult}: \\
        Client pollt Ergebnis per \texttt{task\_id}.
        Dispatcher gibt aktuellen Zustand und Ergebnis zurück.
\end{enumerate}

\begin{figure}[H]
	\centering
	\includegraphics[scale=0.65]{AblaufdiagrammNormal.drawio.png}
	\caption{Erfolgsfall-Ablaufdiagramm}
\end{figure}
